% Размер бумаги А4, шрифт 12 пунктов
\documentclass[a4paper,12pt]{report}
% Перекодировка шрифтов
\usepackage[T2A]{fontenc}
% Кодировка: utf8
\usepackage[utf8]{inputenc}
% Используем русский и английский языки с переносами
\usepackage[english,russian]{babel}
% Подключаем нужные пакеты расширений
\usepackage{amssymb,amsfonts,amsmath,mathtext,cite,enumerate,float}

% Вставка рисунков
\usepackage{graphicx}
% Путь к рисункам
\graphicspath{{../pictures/}}

% Рыба
\usepackage{lipsum}
% Подсчёт объектов
\usepackage{totcount}
\regtotcounter{page}
\regtotcounter{figure}
\regtotcounter{table}

% Изменение размеров шрифтов
\usepackage{scrextend}
% Размер по умолчанию — 14 пунктов, интервал полуторный (14*1.25)
\changefontsizes[17.5pt]{14pt}
% Отступ первой строки параграфа после начала раздела
\usepackage{indentfirst}
% Отступ начала параграфа
\setlength{\parindent}{1cm}

% Заменяем библиографию с квадратных скобок на точку
\makeatletter
\renewcommand{\@biblabel}[1]{#1.}
\makeatother

% Геометрия страницы
\usepackage{geometry}
% Поля
\geometry{left=3cm}
\geometry{right=1cm}
\geometry{top=2cm}
\geometry{bottom=2cm}

% Меняем везде перечисления на цифра.цифра
\renewcommand{\theenumi}{\arabic{enumi}.}
\renewcommand{\labelenumi}{\arabic{enumi}.}
\renewcommand{\theenumii}{\arabic{enumii}.}
\renewcommand{\labelenumii}{\arabic{enumi}.\arabic{enumii}.}
\renewcommand{\theenumiii}{\arabic{enumiii}.}
\renewcommand{\labelenumiii}{\arabic{enumi}.\arabic{enumii}.\arabic{enumiii}.}

% Изменяем \part и \chapter из report.cls чтобы они не выводили нумерованные заголовки
% Рецепт отсюда: http://tex.stackexchange.com/questions/70621/part-and-chapter-commands-output-both-numbered-and-named-headings/70624
\makeatletter
    \def\@part[#1]#2{%
        \ifnum \c@secnumdepth >-2\relax
        \refstepcounter{part}%
        \addcontentsline{toc}{part}{#1}%
        \else
        \addcontentsline{toc}{part}{#1}%
        \fi
        \markboth{}{}%
        {
            \centering
            \interlinepenalty \@M
            \normalfont
            \Huge \bfseries #2\par}%
        \@endpart}
    \def\@chapter[#1]#2{\ifnum \c@secnumdepth >\m@ne
        \refstepcounter{chapter}%
        \typeout{\@chapapp\space\thechapter.}%
        \addcontentsline{toc}{chapter}%
        {#1}%
        \else
        \addcontentsline{toc}{chapter}{#1}%
        \fi
        \chaptermark{#1}%
        \addtocontents{lof}{\protect\addvspace{10\p@}}%
        \addtocontents{lot}{\protect\addvspace{10\p@}}%
        \if@twocolumn
        \@topnewpage[\@makechapterhead{#2}]%
        \else
        \@makechapterhead{#2}%
        \@afterheading
        \fi}
    \def\@makechapterhead#1{%
        \vspace*{50\p@}%
        {\parindent \z@ \raggedright \normalfont
            \interlinepenalty\@M
            \Huge \bfseries #1\par\nobreak
            \vskip 40\p@
    }}
    \renewcommand{\thesection}{\arabic{section}} % Удаляем \thechapter из \thesection
\makeatother

% Для веб-адресов в библиографии
\usepackage{url}
% Библиография в порядке упоминания в тексте
\bibliographystyle{unsrt}

\begin{document}
% Титульный лист
\begin{titlepage}
    \newgeometry{top=2cm,bottom=2cm,left=2cm,right=2cm}
    \newpage
    
    \begin{center}
        Федеральное агенство по образованию Российской Федерации \\
        Московский Государственный Технический Университет \\*
        имени Н.Э.Баумана \\*
        \vspace{-12mm}
        \begin{figure}[h]
            \center{\includegraphics[width=0.2\linewidth]{symbol}}
        \end{figure}
        \vspace{-16mm}
        \hrulefill
    \end{center}
    \center{Факультет «Робототехника и комплексная автоматизация»\\
    		Кафедра «Системы Автоматизированного Проектирования»}
    \begin{center}
        \Large Пояснительная записка к курсовому проекту на тему:
    \end{center}
    
    \vspace{2.5em}
    
    \begin{center}
        \textsc{\textbf{Инструментарий коллективной оптимизации программ.}}
    \end{center}
    
    \vspace{6em}
    
    \begin{flushleft}
        \hspace{8.5cm}Исполнитель \hrulefill Панков М.К. \\
        \vspace{1.5em}
        \hspace{8.5cm}Руководитель \hrulefill Карпенко А.П.\\
    \end{flushleft}
    
    \vspace{\fill}
    
    \begin{center}
        Москва, \\*
        2012
    \end{center}

\end{titlepage}
% Введение
\section{Введение}
\subsection{Мотивация}
Ниже приведены основные причины заинтересованности в данной задаче.
\begin{itemize}
    \item Многие исследования в области применения машинного обучения в компиляции \cite{Dubach2009PCO,Thomson2009RTT,Stephenson2003MOI} показывают большой потенциал и востребованность решения данной задачи, но плохо систематизированы.
    \item Инженерная задача по созданию инфраструктуры для коллективной оптимизации не решена до сих пор \cite{Fursin2010COP}.
    \item Конечным пользователям компиляторов мешает именно отсутствие инструментария, который помог бы им пользоваться обучаемым компилятором так же, как обычным (используя интерфейс «чёрного ящика»).
    \item Производители аппаратного и программного обеспечения, такие как Intel, IBM, ARM заинтересованы в решении такой задачи \cite{IntelExascale}.
\end{itemize}


\subsection{Цель}
Конечная цель, которой хотелось бы достичь --- разработать систему сбора, систематизации, формализации данных о производительности компилируемых программ в зависимости от настроек компилятора, а также выполняющую функции поддержки базы знаний и обучения. Цели данной системы представлены ниже.
\begin{itemize}
    \item Предсказание производительности программ.
    \item Оптимизация производительности программ путём изменения настроек компилятора.
\end{itemize}

На данном, начальном, этапе работы, рассматривается только 1 критерий оптимальности --- производительность, а именно --- полное время выполнения программы.

Однако, сравнивая стоящую перед нами задачу с текущими успехами в области \cite{Fursin2010COP}, обнаруживаем, что достижение данной цели в рамках курсового проекта не представляется возможным. В связи с этим разработка ограничена учебным заданием и продолжение работы будет осуществляться в рамках дипломного проекта.

Итак, цель данной работы --- разработать основные компоненты расширяемого инструментария для проведения экспериментов по оптимизации программ (в ручном и автоматическом режиме) и прозрачного использования данных из экспериментов в промышленной компиляции.


\subsection{Существующие работы в области}
Как было сказано выше, многие исследования в области применения машинного обучения в компиляции \cite{Dubach2009PCO,Thomson2009RTT,Stephenson2003MOI} показывают большой потенциал и востребованность решения данной задачи, но плохо систематизированы. Вследствие этого основной проблемой является реализация общей базы данных программ и экспериментов и каркаса инструментария для проведения исследований. Реализация всех конечных сценариев, необходимых в исследовании, не является приоритетной задачей, но должно быть осуществимо достаточно простым способом силами самих программных инженеров-конечных пользователей системы и других исследователей.

Ниже приведены некоторые отличия данной работы от существующих.
\begin{itemize}
    \item Исследователям и инженерам нужен в первую очередь минимальный функционал для решения описываемой задачи. Графический интерфейс не является необходимым. В сосредоточенности на традиционном текстовом командном интерфейсе основное отличие этой работы от работы \cite{Fursin2010COP}.
    \item Инструментарий не может реализовывать все необходимые сценарии работы с самого начала. Поэтому цель --- сделать систему компактной, модульной и документированной, что обеспечит хорошую поддерживаемость и расширяемость. В этом также одно из отличий от работы \cite{Fursin2010COP}.
    \item Разработка системы сфокусирована на применении локально в изоляции от других исследователей, поскольку это --- первоначально необходимый сценарий использования. Разработка многопользовательских сценариев является следующим шагом. Однако используемая технологическая база позволяет относительно просто расширить инструментарий для использования в распределённых системах.
\end{itemize}


\subsection{Задачи}
Задачи, решаемые в данной работе, перечислены далее в примерном порядке реализации.
\begin{enumerate}
    \item Установить репозиторий исходных кодов программ.
    \item Создать и настроить базу данных для хранения экспериментов.
    \item Реализовать интерактивное текстовое окружение для проведения экспериментов, обработки и сохранения данных.
    \item Реализовать построение графиков для визуального представления данных экспериментов.
\end{enumerate}
% Инструментарий коллективной оптимизации
\section{Инструментарий коллективной оптимизации программ}
% TODO


\subsection{Расширяемость инструментария}
Инструментарий должен обеспечивать расширяемость в смысле возможности добавления новых модулей сборки, запуска и анализа данных. Это должно обеспечиваться описанием взаимодействия различных модулей системы. На данный момент эта задача не решается, поскольку более важным является получение результатов хотя бы на ограниченном множестве программ с фиксированным набором модулей сборки, запуска и анализа. На данный момент введение обобщённой архитектуры является преждевременным.


\subsection{Платформа для использования инструментария}
В идеальном варианте инструментарий должен поддерживать максимальное число различных программно-аппаратных платформ, как настольных, так и мобильных.

На начальном этапе инструментарий поддерживает ОС Linux на x86-совместимом процессоре ввиду распространённости данной аппаратной платформы, а также удобства разработки и ведения исследований под Linux.


\subsection{Эксперимент по оптимизации программы}
Эксперимент состоит в следующем.

\begin{enumerate}
	\item Система собирает программу с определёнными настройками сборки. В случае компилятора gcc они включают в себя \cite{gcc-options}:

	\begin{enumerate}
		\item компилятор (команда для запуска);
		\item базовый уровень оптимизации (флаг '-O[n]', где n --- уровень оптимизации);
		\item набор флагов тонких настроек оптимизации (флаги семейства '-f[name]', где name --- имя определённого набора оптимизаций). Некоторые из этих флагов также имеют числовые параметры;
		\item путь к заголовочным файлам (опция '-I');
		\item путь к исходным файлам;
		\item параметры определения макросов (флаги вида '-D[MACRO]', где MACRO --- имя определяемого макроса);
		\item путь к исполняемому файлу, производимому компилятором.
	\end{enumerate}

	\item Система запускает программу в контролируемом окружении с определёнными настройками запуска и производит измерение интересующих метрик исполнения программы. Они могут включать в себя:

	\begin{enumerate}
		\item полное время исполнения программы;
		\item время исполнения программы по функциям;
		\item полный объём памяти, занимаемой программой и данными;
		\item объём памяти, занимаемой кодом программы;
	\end{enumerate}

	На данный момент нас будет интересовать полное время исполнения.
	Настройки включают в себя:

	\begin{enumerate}
		\item путь к исполняемому файлу;
		\item источник ввода для стандартного потока ввода программы (перенаправление stdin);
		\item потоки вывода для стандартного потока вывода и стандартного потока ошибок программы (перенаправление stdout и stderr соответственно);
		\item аргументы запуска, в т.ч. путь к обрабатываемому файлу данных (например, путь к файлу изображения для программы сжатия изображений).

	\end{enumerate}

	\item Система сохраняет данные о сборке и запуске программы в базу данных для последующего анализа и обработки.
\end{enumerate}


\subsubsection{Время исполнения программы}
Точное измерение времени исполнения программы может представлять сложности ввиду возможных колебаний из-за изменения загрузки системы другими задачами. Принципиально, рекомендуется выполнять эксперименты на системе, не занятой другими задачами. Однако даже в таком случае системные процессы или разница в решениях, принятых планировщиком процессов, могут оказать значительное воздействие на измерение времени.

Для устранения описанных проблем необходимо производить несколько запусков программы с последующей статистической обработкой. Как показывают работы \cite{runs-number}, 5 измерений обычно достаточно для получения достаточной точности.

В случае внешнего измерения времени исполнения влияние оказывают также накладные расходы, связанные с запуском программы из системы. Для их исключения измерять время исполнения лучше изнутри программы и выводить его, например, в стандартный поток ошибок. Таким образом измерение времени производится в пакете тестирования производительности POLYBENCH \cite{polybench}.

Внутреннее измерение времени исполнения обеспечивает большую точность, однако, оно требует поддержки на уровне исходного кода. Очевидно, что для универсального инструментария запуска оптимизационных экспериментов такое решение неприемлемо, поскольку должно быть возможно измерение времени выполнения программ, никак специально не подготовленных для этого.

По этой причине система измеряет время исполнения извне. Невысокая точность измерения при этом устраняется с помощью калибровки (см. ниже). При этом также стоит отметить, что погрешность, вносимая внешним измерением, является систематической и не оказывает влияния на относительное время выполнения разных версий программы, что является наиболее важным в данной задаче.

% FIXME: Не проверено.
% Однако, введение в систему множественных запусков одной и той же программы вносит погрешность другого рода, связанную с т.н. эффектом "прогрева кэша" процессора \cite{cache-heating}. Для устранения этой погрешности между запусками измеряемой программы необходимо запускать некую программу-"заглушку", которая исполняет в достаточно длинном цикле некие (возможно бессмысленные) операции с целью заполнить кэш процессора своими данными. Стоит отметить, что такая программа будет зависеть от конкретной архитектуры ввиду необходимости учитывать размер кэшей процессора. Мы будем придерживаться источника \cite{cache-clearing} относительно необходимости запуска программы, очищающей кэш.

Помимо отмеченных сложностей, есть сложность измерения маленьких интервалов времени. Если время исполнения программы находится на уровне 1 мс (таково временное разрешение системного вызова Unix 'gettimeofday'), то даже если запускать программу несколько раз, колебания измерений окажутся настолько большими, что сделают результат неточным. Для устранения этой проблемы необходимо выполнять запуски программы с накоплением времени до тех пор, пока колебания измеренного времени не достигнут некоторой фиксированной малой величины, определяемой необходимой точностью измерений. Данный процесс называется калибровкой. Схема алгоритма приведена ниже.

\begin{enumerate}
	\item Положить n = 1.
	\item Запустить программу 10^n раз. Вычислить дисперсию времени выполнения. Если дисперсия меньше константы \epsilon и время выполнения многих запусков превысило 10 мсек, сохранить число запусков для последующего использования и закончить работу.
\end{enumerate}



% Название библиографии
\renewcommand\bibname{Список использованных источников}
% Список использованных источников
\bibliography{biblio/ref}
\end{document}